% =====================================================================
% TEXTOS PARA O RELATÓRIO 2 (Overleaf, modelo SBC)
% Arquivo só para copiar e colar: não precisa ir para o repositório.
%
% IMAGENS: suba no Overleaf estes arquivos de resultados/ (mesmos nomes):
%   workload2_frag_externa.png, coalescencia.png, escala.png,
%   barra_heap.png, tempo_malloc.png, min_dados.png
%
% CITAÇÕES: troque as chaves \cite{...} pelas chaves que vocês já usam no
% .bib do relatório 1 (ex.: wilson1995 = [Wilson et al. 1995]).
% As referências novas estão na PARTE C, no fim deste arquivo.
%
% TAMANHO: o roteiro recomenda de 5 a 8 páginas. Se passar, a figura
% barra_heap.png (marcada como OPCIONAL) é a primeira candidata a sair.
% =====================================================================


% =====================================================================
% PARTE A — CORREÇÕES NAS SEÇÕES QUE JÁ EXISTEM NO RELATÓRIO 1
% =====================================================================

% ---------------------------------------------------------------------
% A.1 Introdução: troque a última frase do segundo parágrafo
% ("Aliás, posteriormente, faremos uma avaliação experimental...") por:
% ---------------------------------------------------------------------
Além disso, realizamos uma avaliação experimental das estratégias First-Fit,
Best-Fit, Worst-Fit e Next-Fit, considerando indicadores de fragmentação,
taxa de falha e o custo das operações, com e sem coalescência de blocos.

% ---------------------------------------------------------------------
% A.2 Seção 3.1 (Ambiente): acrescente estes itens à lista
% ---------------------------------------------------------------------
\item Bibliotecas Python: Pandas, NumPy e Matplotlib;
\item Ferramentas de verificação: Valgrind e AddressSanitizer;
\item Flags de compilação: \texttt{-Wall -Wextra -g -std=c11}.

% ---------------------------------------------------------------------
% A.3 Seção 3.3 (Estrutura dos metadados): substitua o código e a lista
% (mantenha o mesmo ambiente de código colorido do relatório 1; aqui está
%  em verbatim só para facilitar a cópia do conteúdo)
% ---------------------------------------------------------------------
\begin{verbatim}
typedef struct bloco {
    size_t tamanho;
    size_t pedido;
    int livre;
    struct bloco* proximo;
    struct bloco* anterior;
} Bloco;
\end{verbatim}

\begin{itemize}
    \item \texttt{size\_t tamanho}: quantidade de bytes de dados do bloco, sem contar o cabeçalho. Ocupa 8 bytes.
    \item \texttt{size\_t pedido}: quantidade de bytes que o usuário realmente pediu. É usado apenas nas estatísticas, pois \texttt{tamanho - pedido} é a fragmentação interna do bloco. Ocupa 8 bytes.
    \item \texttt{int livre}: funciona como uma flag (1 significa livre e 0 significa ocupado). Ocupa 4 bytes, e o compilador acrescenta mais 4 bytes de preenchimento para que o ponteiro seguinte fique alinhado.
    \item \texttt{bloco* proximo}: ponteiro para o próximo bloco. Ocupa 8 bytes.
    \item \texttt{bloco* anterior}: ponteiro para o bloco anterior. Ocupa 8 bytes.
\end{itemize}

% e troque a frase "O tamanho total dos metadados é 32 bytes..." por:
O tamanho total dos metadados é 40 bytes em sistemas de 64 bits. Em relação à
estrutura planejada no relatório inicial, o campo \texttt{padding[4]} foi
substituído pelo campo \texttt{pedido}, necessário para descontar a
fragmentação interna de um bloco no momento em que ele é liberado.

% ---------------------------------------------------------------------
% A.4 Seção 3.7 (Tratamento de erros): acrescente à lista e ao texto
% ---------------------------------------------------------------------
\item \texttt{ERR\_FRAGMENTATION}: há espaço livre suficiente somando todos os blocos livres, mas nenhum bloco contíguo grande o bastante para atender à requisição.

% parágrafo depois da lista:
Cada erro imprime uma mensagem na saída de erro padrão e é registrado na
variável global \texttt{ultimo\_status}, que guarda o resultado da última
chamada de \texttt{meu\_malloc} ou \texttt{meu\_free}. Assim, os testes
automatizados conseguem verificar o tipo de erro sem depender do texto
impresso. A falha por pedido maior que o heap inteiro é tratada como
\texttt{ERR\_OUT\_OF\_MEMORY}.

% ---------------------------------------------------------------------
% A.5 Seção 3.8 (Métricas), parágrafo da taxa de falha: troque por
% ---------------------------------------------------------------------
O total\_solicitacoes é incrementado em toda chamada de alocação com tamanho
maior que zero. Toda alocação que retorna NULL é contada como falha, e as
falhas em que total\_espaco\_livre $\geq$ tamanho\_pedido são contadas também
em separado, como falhas por fragmentação externa. Dessa forma é possível
distinguir a falta real de memória da falha causada apenas pela fragmentação.

% ---------------------------------------------------------------------
% A.6 Seção 3.9 (Organização do projeto): substitua a lista
% ---------------------------------------------------------------------
\begin{itemize}
    \item \texttt{alocador.h} / \texttt{alocador.c}: estrutura do bloco, constantes, funções principais, splitting e coalescência.
    \item \texttt{estrategias.h} / \texttt{estrategias.c}: buscas First-Fit, Next-Fit, Best-Fit e Worst-Fit, escolhidas pela variável global \texttt{estrategia\_atual}.
    \item \texttt{estatisticas.h} / \texttt{estatisticas.c}: contadores das métricas, medição de tempo e gravação das linhas do CSV.
    \item \texttt{testes.c}: testes básicos obrigatórios de funcionamento.
    \item \texttt{driver.c}: lê comandos (\texttt{ALLOC}, \texttt{FREE}, \texttt{DUMP}, \texttt{STATS}) de um arquivo; os cenários de \texttt{exemplos/} são comparados com a saída esperada.
    \item \texttt{workload1.c}, \texttt{workload2.c} e \texttt{workload3.c}: os três padrões de teste descritos na Seção~\ref{sec:experimentos}.
    \item \texttt{referencia\_malloc.c} e \texttt{min\_dados.c}: experimentos complementares (comparação com o \texttt{malloc} real e efeito do tamanho mínimo de divisão).
    \item \texttt{analise.py}: lê os CSVs de \texttt{resultados/} e gera as tabelas e os gráficos.
\end{itemize}

% ---------------------------------------------------------------------
% A.7 Seção 3.10 (Funções principais): acrescente ao código e à lista
% ---------------------------------------------------------------------
void  mem_destroy(void);

\item \texttt{void mem\_destroy(void)}: devolve ao sistema operacional a região reservada por \texttt{mem\_init}. Também é chamada no início de \texttt{mem\_init}, para que uma reinicialização não deixe o heap anterior vazando.

% e acrescente ao fim da seção:
A coalescência pode ser desligada pela variável global
\texttt{coalescencia\_ativa}, o que permite repetir o mesmo workload com e sem
a fusão de blocos.

% ---------------------------------------------------------------------
% A.8 Seção 4 (Conclusão) do relatório 1: apague e use a Seção 6 da PARTE B.
% (corrige também os erros "estrutuda" e "worloads")
% ---------------------------------------------------------------------


% =====================================================================
% PARTE B — SEÇÕES NOVAS (colar depois da Seção 3.10)
% =====================================================================

\subsection{Implementação e decisões de projeto}

A busca por bloco livre foi separada da lógica de alocação: \texttt{meu\_malloc}
alinha o tamanho pedido, chama a função \texttt{buscar}, que executa a
estratégia indicada por \texttt{estrategia\_atual}, e depois realiza o
splitting. As quatro buscas percorrem a mesma lista implícita e contam cada
bloco inspecionado, o que alimenta a métrica de operações por alocação. Além
das três estratégias exigidas, foi implementada a Next-Fit, que recomeça a
busca a partir do bloco em que a busca anterior terminou.

Algumas decisões de projeto merecem destaque:

\begin{itemize}
    \item \textbf{Tamanho mínimo de divisão:} o bloco só é dividido quando a sobra comporta um cabeçalho e pelo menos \texttt{MIN\_DADOS} = 8 bytes de dados. Caso contrário, o bloco é entregue inteiro e a sobra é contabilizada como fragmentação interna.
    \item \textbf{Desempate:} no Best-Fit e no Worst-Fit, quando dois blocos têm o mesmo tamanho, é escolhido o de menor endereço, o primeiro encontrado na lista.
    \item \textbf{Validação de ponteiro:} além de verificar se o endereço está dentro do heap, \texttt{meu\_free} percorre a lista para confirmar que o endereço é o início dos dados de um bloco, com custo $O(n)$.
    \item \textbf{Next-Fit e coalescência:} quando o bloco guardado pela Next-Fit é absorvido por uma fusão, a busca passa a recomeçar do bloco que o absorveu.
    \item \textbf{Consistência:} \texttt{mem\_dump} confere, para cada bloco, se \texttt{proximo->anterior} aponta de volta e se o fim do bloco coincide com o início do próximo, emitindo um alerta caso contrário.
\end{itemize}

\subsection{Validação}

A corretude foi verificada de três formas. Primeiro, o programa
\texttt{testes.c} executa 102 verificações automáticas, cobrindo inicialização,
alinhamento, splitting (inclusive o caso em que a sobra é pequena demais),
coalescência ligada e desligada, a escolha de cada estratégia, heap não
inicializado, falta de memória, falha por fragmentação, ponteiro inválido e
double free. Segundo, o programa \texttt{driver.c} executa doze cenários de
entrada e compara a saída com o resultado esperado, entre eles os cenários do
roteiro com Best-Fit e Worst-Fit, com e sem coalescência. Por fim, testes e
cenários foram executados no Valgrind, sem erros de memória nem vazamentos.

\subsection{Experimentos}\label{sec:experimentos}

Os workloads usam semente fixa (42), de modo que todas as estratégias recebem
exatamente a mesma sequência de pedidos, com tamanhos sorteados entre 10 e 150
bytes:

\begin{itemize}
    \item \textbf{Workload 1:} heap de 4096 bytes e alocações sequenciais, sem liberação, até a primeira falha.
    \item \textbf{Workload 2:} heap de 2000 bytes e 100 operações intercaladas, sendo 60\% alocações e 40\% liberações de blocos ativos sorteados. A fragmentação externa é registrada após cada operação. O mesmo workload é repetido com heaps de 1000, 5000 e 20000 bytes e com 100, 1000 e 10000 operações para o experimento de escala.
    \item \textbf{Workload 3:} vinte blocos de 64 bytes ocupando o heap inteiro; os blocos de índice par são liberados e é feito um pedido de 640 bytes, a soma dos buracos. Depois, os blocos ímpares também são liberados e o pedido é repetido.
\end{itemize}

Os workloads 2 e 3 foram executados com a coalescência ligada e desligada. Os
resultados de cada operação são gravados em \texttt{resultados/comparativo.csv}
e analisados pelo script \texttt{analise.py}.


\section{Resultados Experimentais}

A Tabela~\ref{tab:resumo} resume o estado final de cada workload com a
coalescência ligada. Como os contadores são acumulados ao longo da execução,
os valores correspondem à última operação; a fragmentação externa é a média ao
longo de todas as operações.

\begin{table}[ht]
\centering
\caption{Resumo por workload e estratégia (coalescência ligada).}
\label{tab:resumo}
\begin{tabular}{llrrrrr}
\hline
Workload & Estratégia & Aloc. OK & Falhas & Frag. ext. (\%) & Maior livre & Exam./aloc. \\
\hline
1 & First-Fit & 33 & 1 & 0,0 & 0 & 17,47 \\
1 & Next-Fit  & 33 & 1 & 0,0 & 0 & 2,88 \\
1 & Best-Fit  & 33 & 1 & 0,0 & 0 & 17,47 \\
1 & Worst-Fit & 33 & 1 & 0,0 & 0 & 17,47 \\
\hline
2 & First-Fit & 46 & 22 & 21,4 & 80  & 11,59 \\
2 & Next-Fit  & 47 & 21 & 42,7 & 128 & 7,85 \\
2 & Best-Fit  & 46 & 22 & 11,5 & 80  & 13,40 \\
2 & Worst-Fit & 48 & 20 & 40,9 & 48  & 14,87 \\
\hline
3 & First-Fit & 21 & 1 & 24,2 & 1360 & 10,50 \\
3 & Next-Fit  & 21 & 1 & 24,2 & 1360 & 2,73 \\
3 & Best-Fit  & 21 & 1 & 24,2 & 1360 & 10,50 \\
3 & Worst-Fit & 21 & 1 & 24,2 & 1360 & 10,50 \\
\hline
\end{tabular}
\end{table}

No Workload 1, todas as estratégias realizaram 33 alocações antes da primeira
falha, com 131 bytes de fragmentação interna e 1320 bytes gastos com
cabeçalhos, o equivalente a 32,2\% do heap. Como nenhum bloco é liberado, só
existe um bloco livre, o do fim do heap, e todas as estratégias o escolhem.
A única diferença é o custo: First-Fit, Best-Fit e Worst-Fit percorrem a lista
inteira até o último bloco, enquanto a Next-Fit começa a busca no ponto em que
parou e examina em média 2,88 blocos por alocação.

No Workload 2, a fragmentação externa variou bastante ao longo das operações
(Figura~\ref{fig:w2}). O Best-Fit manteve a menor fragmentação externa média
(11,5\%), seguido pelo First-Fit (21,4\%), enquanto Worst-Fit e Next-Fit
ficaram acima de 40\%. A taxa de falha foi parecida entre as estratégias,
entre 29,4\% e 32,4\% dos 68 pedidos, já que o heap de 2000 bytes é pequeno
para a quantidade de alocações.

\begin{figure}[ht]
\centering
\includegraphics[width=\textwidth]{workload2_frag_externa.png}
\caption{Fragmentação externa ao longo das operações do Workload 2 (coalescência ligada).}
\label{fig:w2}
\end{figure}

No Workload 3, o pedido de 640 bytes falhou para todas as estratégias com
fragmentação externa de 90\%: havia 640 bytes livres, mas o maior bloco livre
tinha apenas 64 bytes. Com a coalescência ligada, a liberação dos blocos
ímpares uniu todo o heap em um único bloco e o segundo pedido foi atendido;
sem coalescência, o segundo pedido também falhou, mesmo com o heap totalmente
livre (Figura~\ref{fig:coal}).

\begin{figure}[ht]
\centering
\includegraphics[width=\textwidth]{coalescencia.png}
\caption{Coalescência ligada e desligada: fragmentação externa média no Workload 2 e falhas no Workload 3.}
\label{fig:coal}
\end{figure}

A Tabela~\ref{tab:coal} mostra o efeito da coalescência na fragmentação
externa média do Workload 2.

\begin{table}[ht]
\centering
\caption{Fragmentação externa média no Workload 2 (\%).}
\label{tab:coal}
\begin{tabular}{lrr}
\hline
Estratégia & Sem coalescência & Com coalescência \\
\hline
First-Fit & 48,1 & 21,4 \\
Next-Fit  & 57,1 & 42,7 \\
Best-Fit  & 40,3 & 11,5 \\
Worst-Fit & 67,1 & 40,9 \\
\hline
\end{tabular}
\end{table}

O experimento de escala (Figura~\ref{fig:escala}) mostra que o custo da busca
cresce com o tamanho do heap e com o número de operações, porque mais blocos
passam a existir na lista. Com heap de 20000 bytes e 10000 operações, o
Worst-Fit examinou em média 180,3 blocos por alocação, o Best-Fit 163,4, o
First-Fit 109,9 e o Next-Fit 73,6. Com heap de 1000 bytes, todas as estratégias
ficaram entre 5 e 8 blocos.

\begin{figure}[ht]
\centering
\includegraphics[width=\textwidth]{escala.png}
\caption{Blocos examinados por alocação em função do tamanho do heap e do número de operações.}
\label{fig:escala}
\end{figure}

% ---- FIGURA OPCIONAL (primeira a sair se o relatório passar de 8 páginas) ----
A Figura~\ref{fig:barra} mostra o heap no fim de cada workload, com os blocos
livres em verde e os ocupados em laranja. No Workload 3 sem coalescência, o
heap termina totalmente livre, porém dividido em vinte blocos de 64 bytes.

\begin{figure}[ht]
\centering
\includegraphics[width=\textwidth]{barra_heap.png}
\caption{Estado final do heap em cada workload e estratégia.}
\label{fig:barra}
\end{figure}
% ---- fim da figura opcional ----

O tempo de cada operação foi medido com \texttt{clock\_gettime(CLOCK\_MONOTONIC\_RAW)}.
Nos workloads, a mediana do tempo de alocação foi de 80 ns no Next-Fit,
131 ns no First-Fit, 150 ns no Worst-Fit e 151 ns no Best-Fit, com percentil 95
entre 170 e 300 ns; a liberação ficou entre 95 e 110 ns de mediana. Na
comparação com o \texttt{malloc} da biblioteca padrão (Figura~\ref{fig:malloc}),
usando um heap de 1 MiB e 2000 operações, o \texttt{malloc} real levou cerca de
50 ns por operação, enquanto o nosso alocador levou de 551 ns (First-Fit) a
1753 ns (Worst-Fit) por alocação. A exceção foi a Next-Fit, com 50 ns, porque
nesse cenário quase sempre encontra espaço logo no bloco em que parou.

\begin{figure}[ht]
\centering
\includegraphics[width=\textwidth]{tempo_malloc.png}
\caption{Tempo por operação (p50 e p95) do nosso alocador e do \texttt{malloc} real.}
\label{fig:malloc}
\end{figure}


\section{Discussão}

\textbf{Qual estratégia apresentou menor fragmentação externa no padrão
adversarial?} No Workload 3, as quatro estratégias tiveram exatamente o mesmo
resultado: 90\% de fragmentação externa no momento do pedido grande. Como todos
os buracos têm o mesmo tamanho (64 bytes) e o pedido não cabe em nenhum, não há
escolha a fazer, e a falha é causada pelo padrão de uso, não pela estratégia. A
comparação entre estratégias aparece no Workload 2, no qual o Best-Fit teve a
menor fragmentação externa média (11,5\%). Isso confirma a teoria apenas em
parte: o Best-Fit de fato deixa sobras pequenas, mas, por escolher o menor
bloco que serve, preserva os blocos grandes. Como a métrica usada é
$1 - \mbox{maior\_livre}/\mbox{total\_livre}$, preservar o maior bloco mantém a
fragmentação baixa. O Worst-Fit, ao contrário, sempre corta o maior bloco e
obteve a maior fragmentação externa média (40,9\%) \cite{wilson1995}.

\textbf{Qual o custo de cada estratégia à medida que o heap fica mais
fragmentado?} Todas as estratégias são $O(n)$ no número de blocos, mas Best-Fit
e Worst-Fit sempre percorrem a lista inteira, enquanto o First-Fit para no
primeiro bloco que serve. Por isso, no experimento de escala, o First-Fit
examinou cerca de 33\% menos blocos que o Best-Fit e 39\% menos que o
Worst-Fit no maior heap. A Next-Fit foi a mais barata por não recomeçar sempre
do início, mas teve fragmentação externa mais alta, pois espalha as alocações
pelo heap.

\textbf{Como a coalescência afeta os resultados?} Sem coalescência, a
fragmentação externa média do Workload 2 aumentou para todas as estratégias,
chegando a 67,1\% no Worst-Fit, e no Workload 3 o pedido final falhou mesmo com
todo o heap livre. A coalescência também muda a escolha das estratégias: no
cenário 5 do roteiro, com Best-Fit sem coalescência, o pedido de 40 bytes usou
o buraco de 56 bytes; com coalescência, esse buraco havia se fundido com o
bloco livre do fim (576 bytes), e o Best-Fit passou a escolher o bloco de 104
bytes. Já no cenário 6, o Worst-Fit escolheu o bloco do fim e deixou uma sobra
grande, de 400 bytes, reaproveitável.

\textbf{Qual o tamanho do cabeçalho e o overhead de metadados?} O cabeçalho tem
40 bytes. No Workload 1, com blocos de 10 a 150 bytes, os cabeçalhos ocuparam
32,2\% do heap. O overhead poderia ser reduzido removendo o campo
\texttt{pedido}, usado apenas nas estatísticas (32 bytes); guardando o estado
livre/ocupado no bit menos significativo de \texttt{tamanho}, que é sempre
múltiplo de 8 (24 bytes); ou usando \textit{boundary tags} no lugar dos
ponteiros \cite{knuth1997}.

\textbf{Qual o valor de tamanho mínimo de divisão e o seu efeito?} Foi escolhido
\texttt{MIN\_DADOS} = 8 bytes, o menor bloco alinhado possível. Para avaliar a
escolha, o padrão do Workload 2 foi executado com 1000 operações e 30 sementes
para os valores 8, 32, 64 e 128 (Figura~\ref{fig:mindados}). Ao passar de 8 para
128 bytes, a fragmentação interna média do Best-Fit subiu de 151 para 443
bytes, porque mais sobras deixam de virar blocos, enquanto a fragmentação
externa caiu de 31,8\% para 14,1\%, já que surgem menos blocos livres pequenos.
A taxa de falha praticamente não mudou (de 32,3\% para 33,0\%). O valor 8 foi
mantido por minimizar a fragmentação interna.

\begin{figure}[ht]
\centering
\includegraphics[width=\textwidth]{min_dados.png}
\caption{Efeito do tamanho mínimo de divisão (média de 30 sementes).}
\label{fig:mindados}
\end{figure}

\textbf{Como o alocador detecta double free e ponteiros inválidos?} O ponteiro
é inválido se estiver fora dos limites do heap ou se, ao percorrer a lista, não
corresponder ao início dos dados de nenhum bloco. O double free é detectado
pela flag \texttt{livre} do cabeçalho. Há limites: a validação custa $O(n)$; o
uso após a liberação (\textit{use after free}) não é detectado; um bloco
liberado que já foi fundido ao anterior é acusado como ponteiro inválido, e não
como double free; e, se um novo bloco for criado exatamente no mesmo endereço,
um ponteiro antigo liberaria o bloco novo sem que o erro fosse percebido.


\section{Conclusão}

Neste trabalho foi implementado um alocador de memória customizado sobre um
heap simulado, com lista duplamente encadeada de blocos, alinhamento de 8
bytes, splitting, coalescência e as estratégias First-Fit, Best-Fit, Worst-Fit e
Next-Fit. A corretude foi verificada por 102 testes automáticos, doze cenários
de entrada e execução no Valgrind.

Os experimentos mostraram que nenhuma estratégia é superior em todos os
aspectos: o Best-Fit obteve a menor fragmentação externa no padrão aleatório,
o First-Fit teve custo menor que Best-Fit e Worst-Fit por parar cedo, e a
Next-Fit foi a mais rápida, porém com maior fragmentação. A coalescência se
mostrou o fator de maior impacto, reduzindo a fragmentação externa em todas as
estratégias e sendo indispensável no padrão adversarial.

Entre as principais dificuldades estiveram a manutenção correta dos quatro
ponteiros no splitting e na fusão de blocos e a validação de ponteiros
inválidos. Como limitações, o heap tem tamanho fixo, a validação de ponteiros
é linear e o cabeçalho de 40 bytes representa um overhead alto para blocos
pequenos, o que poderia ser reduzido com as técnicas discutidas na Seção 5.

% ---- OPCIONAL: declaração de uso de IA (o professor autorizou) ----
% Ferramentas de inteligência artificial foram utilizadas, com autorização do
% professor, como apoio na revisão do código e na redação deste relatório.


% =====================================================================
% PARTE C — REFERÊNCIAS NOVAS (colar no arquivo .bib)
% wilson1995 já existe no relatório 1: só confira a chave.
% =====================================================================
%
% @book{knuth1997,
%   author    = {Knuth, Donald E.},
%   title     = {The Art of Computer Programming, Volume 1: Fundamental Algorithms},
%   edition   = {3},
%   publisher = {Addison-Wesley},
%   year      = {1997}
% }
